%%%%%%%%%%%%%%%%%%%%%%%%%%%%%%%%%%%% Chapter INTRODUÇÃO

\chapter{Introdução} 	% Main chapter title
\label{Chapter1} 		% For referencing the chapter elsewhere, usage \ref{Chapter1}

A este capítulo introdutório  a contextualização e organização do documento descrevem o relatório de uma forma geral, enquanto que os objetivos e o planeamento demonstram mais ao pormenor o estudo e a análise que foi  feito ao longo deste relatório.

%%%%%%%%%%%%%%%%%%%%%%%%%%%%%%%%%%%% Contextualização

\section{Contextualização}
\label{sec:Ch1.1}

\noindent No âmbito da \ac{uc} de \ac{rinte},  será realizado o trabalho nº1, sobre uma tecnologia de suporte ao desenvolvimento de aplicações.~A elaboração deste  relatório técnico-científico é sobre a tecnologia \textit{JavaFX Framework}.~A entrega deste relatório será acompanhada também por código e apresentação exclusivamente para esta \ac{uc} de \ac{rinte} da autoria própria do aluno, que realizou o trabalho, Eduardo Miguel Moreira Junqueira.

\noindent \textit{JavaFX}, uma plataforma de software criada por Sun Microsystems em 2008. Em 2011 tornou-se \textit{open-source} e no ano seguinte migrou para Oracle Corporation's \ac{jdk}. Em 2018 Oracle Corporation's, decidiu não agregar mais esta tecnologia e então formou-se a comunidade representante atualmente com a parceria da GLUON, nomeada de \textit{OpenJFX}\cite{Pereira2020},\cite{OpenJFX}.%JavaFX. (n.d.). Community OpenJFX. Retrieved October 29, 2026, from https://openjfx.io/



%%%%%%%%%%%%%%%%%%%%%%%%%%%%%%%%%%%% Objetivos

\section{Objetivos}
\label{sec:Ch1.2}
 %apresentar de forma clara os objetivos que se propõe atingir
\noindent O objetivo principal deste trabalho prático é estudar e  implementar a tecnologia \textit{JavaFX Framework} de suporte ao desenvolvimento de aplicações distribuídas.

\noindent Esta aplicação distribuída, que será criada como exemplo demonstrativo da tecnologia em estudo, tem como âmbito seguir uma arquitetura de três camadas, ou seja, \textbf{I-} cliente gráfico em \textit{JavaFX},  \textbf{II-} servidor em \textit{Tomcat},  \textbf{III-} base de dados em \textit{MySQL}.

\noindent Inicialmente será realizado um teste de implementação em conjunto com a pesquisa, para verificar a tecnologia a funcionar com as devidas bibliotecas necessárias.

\noindent De seguida, será implementado um projeto nomeado "Gestão de Tarefas",~com dois tipos de clientes, o administrador e o utilizador.

\begin {enumerate}
\item Administrador:~usa métodos \ac{http} do tipo \textit{GET,POST,PUT,DELETE}  para controlar todas as tarefas.
Também usa métodos \ac{http} do tipo \textit{GET,POST,DELETE} para controlar todos os utilizadores.
\item Utilizador:~usa métodos \ac{http} do tipo \textit{GET,POST,PUT}  para controlar apenas as suas tarefas.
\end {enumerate}
%%%%%%%%%%%%%%%%%%%%%%%%%%%%%%%%%%%% Planeamento

\section{Planeamento}
\label{sec:Ch1.3}

A tabela seguinte \ref{tab:cronograma}, demonstra as principais tarefas que foram realizadas de forma autónoma, ao longo de 5 semanas, após a atribuição da tecnologia em estudo.
% Cores próximas das usadas no Excel
\definecolor{fundo}{HTML}{C9C9C9}
\definecolor{pesquisa}{HTML}{AA55A0}
\definecolor{relatorio}{HTML}{32A34A}
\definecolor{implementacao}{HTML}{1084A7}
\definecolor{validacao}{HTML}{B75A28}

% Célula ativa: recebe uma cor e apresenta o valor 1
\newcommand{\ativo}[1]{\cellcolor{#1}1}


\begin{table}[htbp]
    \centering
    \renewcommand{\arraystretch}{1.5}
    \setlength{\tabcolsep}{8pt}

\begin{tabular}{|>{\columncolor{fundo}}l|*{5}{>{\columncolor{fundo}}c|}}
        \hline
        Tarefas
        & 28/9 a 4/10
        & 5/10 a 11/10
        & 12/10 a 18/10
        & 19/10 a 25/10
        & 26/10 a 1/11 \\
        \hline

        Pesquisa
        & \ativo{pesquisa}
        & \ativo{pesquisa}
        & \ativo{pesquisa}
        & \ativo{pesquisa}
        & \ativo{pesquisa} \\
        \hline

        Relatório
        & 0
        & 0
        & 0
        & \ativo{relatorio}
        & \ativo{relatorio} \\
        \hline

        Implementação
        & 0
        & 0
        & \ativo{implementacao}
        & \ativo{implementacao}
        & \ativo{implementacao} \\
        \hline

        Validação
        & 0
        & 0
        & 0
        & 0
        & \ativo{validacao} \\
        \hline
    \end{tabular}

    \caption{Cronograma das tarefas.}
    \label{tab:cronograma}
\end{table}


%%%%%%%%%%%%%%%%%%%%%%%%%%%%%%%%%%%% SECTION 5

\section{Organização do Documento}
\label{sec:Ch1.4}

\noindent Nos seguintes capítulos, o problema contém a definição e os requisitos quer funcionais, quer técnicos. A solução explicita as ferramentas e tecnologias utilizadas, identificando a arquitetura, os modelo de dados, \textit{front-end, back-end}. Os testes e avaliação demonstram as funcionais o desempenho e também a usabilidade. A conclusão retrata o sumário dos resultados com as devidas limitações e o desenvolvimento futuro do trabalho nº2.	




